\documentclass[10pt,a4paper]{amsart}
\usepackage[margin=19mm]{geometry}
\usepackage{amsmath,amssymb,mathtools}
\usepackage[scr=rsfs,cal=euler]{mathalfa}
\usepackage{xfrac}
\usepackage[unicode,hidelinks,bookmarks=false]{hyperref}
\newcommand{\ie}{i.e.\ }
\newcommand{\cf}{cf.\ }
\newcommand{\resp}{respectively\ }
\newcommand{\Subsection}{\subsection}
\providecommand{\nobreakdash}{\nobreak\hbox{-}}
\setlength{\emergencystretch}{2em}
\multlinegap0pt
\usepackage{amssymb}
\usepackage[normalem]{ulem}
\usepackage[shortlabels]{enumitem}
\newtheorem{assumption}{Assumption}
\newtheorem{lemma}{Lemma}
\title{Splitting schemes and estimators for stochastic differential equations with H\"older multiplicative noise}
\author{Bowen Fang, Dario Span\`o, Massimiliano Tamborrino}
\date{Source: arXiv:2605.16900v2 (CC BY 4.0)}
\begin{document}
\maketitle

\noindent\emph{Source and licence.} Splitting schemes and estimators for stochastic differential equations with H\"older multiplicative noise. Version arXiv:2605.16900v2 (CC BY 4.0), distributed by arXiv under the Creative Commons Attribution 4.0 International licence, http://creativecommons.org/licenses/by/4.0/, as recorded in the arXiv licence metadata for this exact version and shown on the abstract page https://arxiv.org/abs/2605.16900v2. Full licence notice: \texttt{licenses/arxiv-2605.16900-CC-BY-4.0.txt}.\par\smallskip

\noindent\textit{Editorial scope.} Counted unit: the article's lemma lemma:LT\_onestep\_stop (source0.tex lines 303--313). The article's complete original proof of it is delivered with it (source0.tex lines 741--763, one \texttt{proof} environment of 23 lines, which the article places away from the statement). No mathematics is written, repaired or re-derived: the statement and the proof are the article's own lines, in the article's own notation, and no retained line is substituted or re-flowed.  Retained uncounted as the source-local context the proof needs: lines 59--61; lines 134--148; lines 667--672; lines 677--679.  Nothing was omitted from any retained span, so the package carries no omission record.  No retained line contains a citation, so the wrapper carries no bibliography and no bibliography entry.  No retained line contains a figure environment, an image-inclusion command or drawing code, so no image is delivered and the package declares no asset dependency.  Editorial formatting only, in addition: an AMS \texttt{amsart} preamble that restates the article's own declarations and macros used by the retained lines, namely the environments \texttt{assumption}, \texttt{lemma} on separate counters, together with the standard packages those lines need.  The retained environments print in this document as Assumption 1, Lemma 1 in order of appearance, whereas the article prints the same items with the numbering of its own full document.  The complete native source of the article as distributed in the arXiv e-print for this exact version is bundled unchanged as \texttt{sources/200752/source0.tex}.
\begin{equation}\label{eqn:the sde}
    dX_t=f(X_t;\theta)dt+g(X_t;\sigma)dW_t,
\end{equation}

\begin{assumption}\label{assump:main}
    For SDE \eqref{eqn:the sde}, the functions $f,g:\mathbb{R}\to \mathbb{R}$ satisfy  
    \begin{itemize}
        \item[(i)] $f(x)=f_1(x)+f_2(x)$, where $f_1,f_2$ are continuous locally Lipschitz functions, i.e.,
        for any $R>0$ and $\max(|x_1|,|x_2|)\leq R$, there exists an arbitrary constant $C_R$ that only depends on $R$ such that
        \begin{equation}
            \left|f_i(x_1)-f_i(x_2)\right|\leq C_R\left(|x_1-x_2|\right), \quad i\in \{1,2\},\label{eqn:f condition}
        \end{equation}
        \item[(ii)] The function $g$ is locally H\"older with exponent $q\in(0,1/2)$, i.e., for any $R>0$ and $\max(|x_1|,|x_2|)\leq R$, there exists $q\in (0,\frac{1}{2})$ such that
        \begin{equation}
            \left|g(x_1)-g(x_2)\right|\leq C_R \left(|x_1-x_2|+|x_1-x_2|^q\right)\label{eqn:g condition}.
        \end{equation}
    \end{itemize}
    
    \end{assumption}

\begin{lemma}\label{lemma:open_p_bracket}
    Suppose $\left(\mathbb{B},\|\cdot\|_{\mathbb{B}}\right)$ is a Banach space with respect to the norm $\|\cdot\|_{\mathbb{B}}$. Let $X_i\in \mathbb{B},\; i\in \{1,\dots,n\}$. Then, for $n\in \mathbb{N},p\geq 1$, 
    \begin{equation}\label{equation:open pth bracket}
        \left\|\sum_{i=1}^n X_i\right\|_{\mathbb{B}}^p \leq n^{p-1} \sum_{i=1}^n\left\|X_i\right\|_{\mathbb{B}}^p.
    \end{equation}
\end{lemma}

    \begin{equation}\label{eqn:BDG inequality}
        c_p \mathbb{E}[|y(t)|^{\frac{p}{2}}] \leq \mathbb{E}\left(\sup _{0 \leq s \leq t}|x(s)|^p\right) \leq C_p \mathbb{E}[|y(t)|^{\frac{p}{2}}].
    \end{equation}

\begin{lemma}\label{lemma:LT_onestep_stop}
     Denote $x\wedge y=\min(x,y),  \hat{t}=t_k$ for any $ t\in [t_k,t_{k+1}), k=0,\ldots, n-1$, and define $\theta_R=\inf \left\{t \in[0, T]:\left|X^{LT}_t\right|>R\right.$ or $\left.\left|X^{LT}_{\hat{t}}\right|>R\right\}$, for $R>0$. If Assumption \ref{assump:main} holds, then, for any $p>0$ and $t\in [0,T]$, we have
    \begin{equation}\label{eqn:LT onestep bound}
\mathbb{E}\left[\left|X^{LT}_{t \wedge \theta_R}-X^{LT}_{\widehat{t \wedge \theta_R}}\right|^p\right] \leq C_{R,p} h^{p / 2},
 \end{equation}
which  implies 
\begin{equation}\label{eqn:LT onestep convergence rate}
    \lim_{h\rightarrow 0}\sup _{t \in\left[t_{k}, t_{k+1}\right]} \mathbb{E}\left[\left|X^{LT}_{t \wedge \theta_R}-X^{LT}_{\widehat{t \wedge \theta_R}}\right|^p\right]=O\left(h^{p / 2}\right), 
\end{equation}
where $C_{R,p}$ is a constant depending only on $R$ and $p$, and $O(h)$ denotes functions of order $h$.
\end{lemma}

\begin{proof}[Proof of Lemma \ref{lemma:LT_onestep_stop}]
If $t\leq \theta_R$, i.e. $t\wedge\theta_R=t$,  by Assumption \ref{assump:main}(i) on $f_1(x)$, we have
    \begin{equation*}
        \mathcal{R}(h^2;X^{\textrm{LT}}_{t\wedge \theta_R})\leq C_Rh^2.
    \end{equation*}
If $t>\theta_R$, i.e. $t\wedge \theta_R=\theta_R\in [t_k,t_{k+1})$ for some $k\in \{0,1,\dots,N-1\}$, and $p\geq1$, we have
    \begin{align*}
         \left|X^{\textrm{LT}}_{t \wedge \theta_R}-X^{\textrm{LT}}_{\widehat{t \wedge \theta_R}}\right|^p \leq &\left|\int_{t_k}^{t\wedge \theta_R} \left[f_1(X^{\textrm{LT}}_{\hat{s}})+f_2(X^{\textrm{LT}}_s)+\frac{1}{h}\mathcal{R}(h^2;X^{\textrm{LT}}_{\hat{s}})\right]ds+\int_{t}^{t\wedge \theta_R}g(X^{\textrm{LT}}_s)dW_s\right|^p\\
         \leq& 4^{p-1}\left(C_R h^p+\left|\int_{t_k}^{t\wedge \theta_R}f_2(X^{\textrm{LT}}_s)ds\right|^p+C_Rh^{2p}+\left|\int_{t}^{t\wedge \theta_R}g(X^{\textrm{LT}}_s)dW_s\right|^p\right)\\
         \leq&C_{R,p}h^p+4^{p-1}h^{p-1}\int_{t_k}^{t\wedge \theta_R}\left|f_2(X^{\textrm{LT}}_s)\right|^pds+C_Rh^{2p}+4^{p-1}\left|\int_{t}^{t\wedge \theta_R}g(X^{\textrm{LT}}_s)dW_s\right|^p\\
         \leq&C_{R,p}h^p+C_Rh^p+C_Rh^{2p}+C_p\left|\int_{t}^{t\wedge \theta_R}g(X^{\textrm{LT}}_s)dW_s\right|^p,
    \end{align*}
    where we use Lemma \ref{lemma:open_p_bracket} in the second inequality, H\"older's inequality (on $f_1$ and $f_2$) in the third step, and Assumption \ref{assump:main}(i) on $f_2(x)$ using the fact that $t\wedge\theta_r\leq t_{k+1}$ and $f_2(x)$ is bounded in $[t_k,t_{k+1})$ for $|x|\leq R$. Now, if we take the expectation of the last term and use BDG's inequality \eqref{eqn:BDG inequality} on $x(s)=\int_t^{t\wedge\theta_R}g(s)dW_s$, we have 
    \begin{equation*}
        \mathbb{E}\left[C_p\left|\int_{t}^{t\wedge \theta_R}g(X^{\textrm{LT}}_s)dW_s\right|^p\right]\leq C_p\mathbb{E}\left|\int_{t_k}^{t\wedge \theta_R}|g(X^{\textrm{LT}}_s)|^{ 2}ds\right|^{\frac{p}{2}}\leq C_{R,p}h^{\frac{p}{2}}.
    \end{equation*}
    Combining, we get \eqref{eqn:LT onestep bound}, and \eqref{eqn:LT onestep convergence rate} follows immediately.\\
    Finally, if $p\in (0,1)$, using Jensen's inequality we have
    \begin{equation}
        \mathbb{E}\left|X^{\textrm{LT}}_{t\wedge \theta_R}-X^{\textrm{LT}}_{\widehat{t\wedge \theta_R}}\right|^p\leq  \left(\mathbb{E}\left|X^{\textrm{LT}}_{t\wedge \theta_R}-X^{\textrm{LT}}_{\widehat{t\wedge \theta_R}}\right|^2\right)^{\frac{p}{2}}\leq C_{R,p}h^{\frac{p}{2}},
    \end{equation}
    which concludes the proof.
\end{proof}

\end{document}
